\subsection{Statement of main result}

In this paper we always use $V$ to denote a finite-dimensional $\F_p$-vector space and we write $\omega=e^{2\pi i/p}$. Given a function $f\colon V\to\C$ and a shift $h\in V$, we write $\partial_h f(x)=f(x+h)\ol{f(x)}$ for the multiplicative derivative.

\begin{defn}
Given a function $f\colon V\to\C$ and $d\geq 2$, the {\textbf{Gowers uniformity norm}} $\|f\|_{U^d}$ is defined by \[\|f\|_{U^d}^{2^d}=\E_{x,h_1,\ldots,h_d\in V}(\partial_{h_1}\cdots\partial_{h_d}f)(x).\]
\end{defn}

This definition has many useful properties, including being a well-defined norm for all $d\geq 2$. Another useful property is the inductive formula $\|f\|_{U^d}^{2^d}=\E_h \|\partial_h f\|_{U^{d-1}}^{2^{d-1}}$.
See \cite[Lemma B.1]{TZ12} for more properties and references about the Gowers uniformity norms.

Consider a function $f\colon V\to G$ where $G$ is an abelian group. Given a shift $h\in V$, we write $\Delta_hf(x)=f(x+h)-f(x)$ for the additive derivative. It follows from the definition of the Gowers uniformity norm that a 1-bounded function $f\colon V\to\C$ satisfies $\|f\|_{U^d}\leq 1$ with equality if and only if $f=e^{2\pi i P}$ where $P\colon V\to\R/\Z$ satisfies $\Delta_{h_1}\Delta_{h_2}\cdots \Delta_{h_d}P(x)=0$. 

\begin{defn}
A \textbf{non-classical polynomial} of degree at most $k$ is a map $P\colon V\to\R/\Z$ that satisfies \[(\Delta_{h_1}\cdots \Delta_{h_{k+1}}P)(x)= 0\] for all $h_1,\ldots,h_{k+1},x\in V$.
We write $\Poly_{\leqslant k}(V\to\R/\Z)$ for the set of non-classical polynomials of degree at most $k$.
\end{defn}

A classical polynomial is a map $P\colon V\to\F_p$ satisfying the same definition. Thus we use $\Poly_{\leqslant k}(V\to\F_p)$ to denote the set of classical polynomials of degree at most $k$. In this paper we commonly abuse notation to identify $\F_p$ with $\{0,1/p,\ldots,(p-1)/p\}\subset\R/\Z$ when convenient. In this way we consider $\Poly_{\leqslant k}(V\to\F_p)$ to be a subset of $\Poly_{\leqslant k}(V\to\R/\Z)$.

\begin{thm}
\label{thm:inverse}
Fix a prime $p$ and $\delta>0$. There exists an $\epsilon>0$ satisfying \[1/\epsilon=\exp\paren{\exp\paren{\exp\paren{O_p\paren{\log(1/\delta)^{O(1)}}}}}\]
such that the following holds. Let $V$ be a finite-dimensional $\F_p$-vector space. Given a function $f\colon V\to\C$ satisfying $\|f\|_\infty\leq1$ and $\|f\|_{U^{4}}>\delta$, there exists a non-classical cubic polynomial $P\in\Poly_{\leqslant 3}(V\to\R/\Z)$ such that \[\abs{\E_{x\in V} f(x)e^{-2\pi i P(x)}}\geq \epsilon.\]
Furthermore, if $p\geq 3$, the polynomial can be taken to be classical.
\end{thm}

This theorem was proved with ineffective bounds by Tao and Ziegler \cite[Theorem 1.10]{TZ12} and with effective bounds for $p\geq 5$ by Gowers and Mili\'cevi\'c \cite[Theorem 7]{GM20}. The bounds we prove are of the same shape as Gowers and Mili\'cevi\'c's result.
